\documentclass[11pt]{article}
\usepackage[a4paper,margin=0.9in]{geometry}
\usepackage{amsmath,amssymb,amsthm,mathtools}
\usepackage[dvipsnames]{xcolor}
\usepackage{enumitem}
\usepackage{booktabs}
\usepackage{tcolorbox}
\usepackage{fancyhdr}
\usepackage{listings}
\usepackage[colorlinks=true,linkcolor=Blue]{hyperref}
\tcbuselibrary{skins,breakable}

\setlist{itemsep=3pt,topsep=4pt}
\pagestyle{fancy}\fancyhf{}
\lhead{\small Newton--Raphson for the MLE}
\rhead{\small Parametric Inference, B.Stat.\ 3rd yr}
\cfoot{\small\thepage}

\newcommand{\E}{\mathbb{E}}
\newcommand{\Var}{\operatorname{Var}}
\newcommand{\Prob}{\mathbb{P}}
\newcommand{\R}{\mathbb{R}}
\newcommand{\Xbar}{\overline{X}}
\newcommand{\iid}{\overset{\text{iid}}{\sim}}
\newcommand{\I}{\mathbb{I}}

\newtcolorbox{idea}[1]{colback=Blue!4,colframe=Blue!55,title={#1},
  fonttitle=\bfseries,breakable}
\newtcolorbox{warn}[1]{colback=BrickRed!5,colframe=BrickRed!60,title={#1},
  fonttitle=\bfseries,breakable}
\newtcolorbox{recipe}[1]{colback=ForestGreen!5,colframe=ForestGreen!55,title={#1},
  fonttitle=\bfseries,breakable}

\lstset{basicstyle=\ttfamily\footnotesize,frame=single,breaklines=true,
  backgroundcolor=\color{black!3},columns=fullflexible}

\begin{document}

\begin{center}
{\LARGE\bfseries Newton--Raphson for the MLE}\\[3pt]
{\large Where it appears in the course, and how to do it}\\[5pt]
{\small Parametric Inference, B.Stat.\ 3rd Year --- Lectures 7, 8, 10 and Homework 2 Q4}
\end{center}

\vspace{4pt}\hrule\vspace{10pt}

\tableofcontents

%%%%%%%%%%%%%%%%%%%%%%%%%%%%%%%%%%%%%%%%%%%%%%%%%%%%%%%%%
\section{Where this appears in the course}

Narrower than it looks. Three mentions in the notes, and exactly \textbf{one} place where you
are actually required to compute anything.

\begin{center}
\renewcommand{\arraystretch}{1.4}
\begin{tabular}{@{}p{3.4cm}p{7.8cm}p{4.0cm}@{}}
\toprule
\textbf{Where} & \textbf{What it says} & \textbf{Status} \\
\midrule
\textbf{Lecture 8} (21 Aug), p.1 &
``Method of Scoring for computing MLE'': the ML equation, the Newton--Raphson iteration, and
replacing the second derivative by its expectation (the Fisher information) &
The only place the \emph{method} is taught \\
\textbf{Lecture 7} (sheet dated 10 Aug), p.3 &
Cauchy: ``what is ML eq? We get a poly eq with no closed sol.'' &
Motivation only \\
\textbf{Lecture 10} (28 Aug), p.1 &
``For any problem involving Cauchy distribution, starting with sample median is good'' &
Starting-value advice \\
\textbf{Homework 2, Q4} &
Generate Cauchy$(0)$ data for $n=10,15,20$ and compute the MLE &
\textbf{The only place you must run it} \\
\bottomrule
\end{tabular}
\end{center}

\begin{idea}{It is never set as a \# class exercise}
All 34 of the professor's ``\#'' exercises are catalogued in
\texttt{PI\_Class\_Exercises\_Solved.pdf}. The Lecture 8 ones are: solve PS1 Q3 with
$\sigma^2$ unknown; derive the Bayes estimate $\frac{p+X}{p+q+n}$; show
$\max_\theta R(\theta,X/n)=\frac{1}{4n}$; complete the constant-risk minimax proof; and the
improper-prior equalizer question. \textbf{None involves Newton--Raphson.}

So: theory in Lecture 8, practice in Homework 2 Q4, nothing else. The full solution to that
homework question, with code and output, is in \texttt{PI\_Assignment\_Solutions.pdf} \S2.
\end{idea}

%%%%%%%%%%%%%%%%%%%%%%%%%%%%%%%%%%%%%%%%%%%%%%%%%%%%%%%%%
\section{Why the Cauchy needs it, and nothing else does}

Every other family in the syllabus has an MLE you can write down: $\Xbar$ for the Normal,
Poisson and Exponential; $X/n$ for the Binomial; $X_{(n)}$ for $U[0,\theta]$. The Cauchy does
not.

With $f(x\mid\theta)=\dfrac{1}{\pi\{1+(x-\theta)^2\}}$,
\[
\ell(\theta)=-\sum_{i=1}^n\log\big\{1+(x_i-\theta)^2\big\}+\text{const},
\]
\[
\boxed{\ \ell'(\theta)=\sum_{i=1}^n\frac{2(x_i-\theta)}{1+(x_i-\theta)^2},
\qquad
\ell''(\theta)=\sum_{i=1}^n\frac{2\big[(x_i-\theta)^2-1\big]}{\big[1+(x_i-\theta)^2\big]^2}. }
\]

\begin{idea}{Why there is no closed form}
Clear denominators in $\ell'(\theta)=0$:
\[
\sum_{i=1}^{n}2(x_i-\theta)\prod_{j\ne i}\big\{1+(x_j-\theta)^2\big\}=0 .
\]
Each summand is a degree-1 factor times $n-1$ quadratics, so this is a polynomial equation of
degree $\boldsymbol{2n-1}$ in $\theta$. For $n\ge3$ there is no general solution by radicals,
and worse, the equation can have \textbf{several roots} --- so ``the MLE'' means the
\emph{global} maximiser, not just any stationary point. Both facts are why the course
introduces a numerical method at all.
\end{idea}

%%%%%%%%%%%%%%%%%%%%%%%%%%%%%%%%%%%%%%%%%%%%%%%%%%%%%%%%%
\section{The two iterations}

\begin{idea}{Newton--Raphson}
Apply Newton's root-finding method to $\ell'(\theta)=0$:
\[
\theta^{(m+1)}=\theta^{(m)}-\frac{\ell'(\theta^{(m)})}{\ell''(\theta^{(m)})}.
\]
Converges \emph{quadratically} near a root --- typically 3--4 iterations to machine precision.
\end{idea}

\begin{idea}{Method of scoring (Lecture 8)}
Replace the observed second derivative by its expectation. Since
$\E_\theta[\ell''(\theta)]=-n\,I(\theta)$,
\[
\theta^{(m+1)}=\theta^{(m)}+\frac{\ell'(\theta^{(m)})}{n\,I(\theta^{(m)})}.
\]
\end{idea}

The professor's stated reason for the substitution is that $I(\theta)$ is often free of
$\theta$ --- in particular \textbf{whenever $\theta$ is a location parameter} --- which makes
the iteration cheaper. But there is a second, more important reason, and it is the point of
\S6: \emph{$-nI(\theta)$ is always negative, so every scoring step is guaranteed to point
uphill.} Newton--Raphson carries no such guarantee.

\subsection{Fisher information for the Cauchy}

With $u=x-\theta$, the score is $\dfrac{\partial}{\partial\theta}\log f=\dfrac{2u}{1+u^2}$, so
\[
I(\theta)=\E\left[\frac{4U^2}{(1+U^2)^2}\right]
=\frac{4}{\pi}\int_{-\infty}^{\infty}\frac{u^{2}}{(1+u^{2})^{3}}\,du .
\]
Substituting $u=\tan t$ (so $du=\sec^2t\,dt$ and $1+u^2=\sec^2t$),
\[
\int_{-\infty}^{\infty}\frac{u^2}{(1+u^2)^3}du
=\int_{-\pi/2}^{\pi/2}\frac{\tan^2 t}{\sec^4 t}\,dt
=\int_{-\pi/2}^{\pi/2}\sin^2t\cos^2t\,dt
=\frac14\int_{-\pi/2}^{\pi/2}\sin^2(2t)\,dt=\frac{\pi}{8},
\]
so
\[
\boxed{\ I(\theta)=\frac4\pi\cdot\frac\pi8=\frac12\ }
\]
\emph{a constant}, free of $\theta$ --- exactly the location-parameter phenomenon Lecture 8
mentions. (Monte Carlo check with $4\times10^6$ draws: $0.4999$.) The scoring iteration
therefore collapses to something you can run on paper:
\[
\boxed{\ \theta^{(m+1)}=\theta^{(m)}+\frac{2}{n}\,\ell'(\theta^{(m)})\ }
\]
with no second derivative anywhere.

\subsection{Why you start at the sample median}

Lecture 10 says to start from $\tilde X_n$. The reason is that the obvious alternative is
useless: the Cauchy is \emph{stable}, so $\Xbar\sim\mathrm{Cauchy}(\theta,1)$ for
\emph{every} $n$ --- averaging does not concentrate at all, and $\Xbar$ is not even consistent.
The sample median \emph{is} consistent (Lecture 10), so it is a starting value that is already
in the right neighbourhood.

%%%%%%%%%%%%%%%%%%%%%%%%%%%%%%%%%%%%%%%%%%%%%%%%%%%%%%%%%
\section{Worked example}

Take $n=5$ with
\[
x=(-2.1,\ -0.6,\ 0.3,\ 1.2,\ 4.8),
\qquad
\text{median}=0.300,\qquad \Xbar=0.720 .
\]
The mean is already dragged off by the single observation $4.8$ --- a preview of why $\Xbar$
is the wrong tool here. Start both iterations at $\theta^{(0)}=0.300$.

\begin{center}
\textbf{Newton--Raphson}\quad $\theta^{(m+1)}=\theta^{(m)}-\ell'/\ell''$

\renewcommand{\arraystretch}{1.3}
\begin{tabular}{@{}ccccc@{}}
\toprule
$m$ & $\theta^{(m)}$ & $\ell'(\theta^{(m)})$ & $\ell''(\theta^{(m)})$ & step \\
\midrule
0 & $0.300000$ & $-0.286530$ & $-1.938398$ & $-0.147818$ \\
1 & $0.152182$ & $-0.003068$ & $-1.881721$ & $-0.001630$ \\
2 & $0.150552$ & $-0.000001$ & $-1.880643$ & $-0.000001$ \\
3 & $0.150551$ & $\approx0$ & $-1.880642$ & --- \\
\bottomrule
\end{tabular}
\end{center}

\[
\boxed{\ \hat\theta=0.1506\ }
\]
Three iterations. Note the quadratic convergence: the step shrinks from $10^{-1}$ to
$10^{-3}$ to $10^{-6}$. A grid search over $[-6,8]$ confirms the global maximiser is
$0.15056$.

\begin{center}
\textbf{Method of scoring}\quad $\theta^{(m+1)}=\theta^{(m)}+0.4\,\ell'(\theta^{(m)})$

\renewcommand{\arraystretch}{1.25}
\begin{tabular}{@{}cccccccc@{}}
\toprule
$m$ & 0 & 1 & 2 & 3 & 4 & 5 & 6 \\
\midrule
$\theta^{(m)}$ & $0.3000$ & $0.1854$ & $0.1590$ & $0.1526$ & $0.1511$ & $0.1507$ & $0.1506$ \\
\bottomrule
\end{tabular}
\end{center}

Same answer, but convergence is \emph{linear} --- the error falls by a roughly constant factor
each step rather than squaring. That is the trade: scoring is slower but needs no second
derivative.

\begin{center}
\renewcommand{\arraystretch}{1.3}
\begin{tabular}{@{}lccc@{}}
\toprule
& sample mean & sample median & MLE \\
\midrule
estimate of $\theta$ & $0.720$ & $0.300$ & $\mathbf{0.151}$ \\
\bottomrule
\end{tabular}
\end{center}

%%%%%%%%%%%%%%%%%%%%%%%%%%%%%%%%%%%%%%%%%%%%%%%%%%%%%%%%%
\section{Code}

\begin{lstlisting}[language=Python]
import numpy as np
rng = np.random.default_rng(2026)

def score(t, x):  return np.sum(2*(x-t)/(1+(x-t)**2))
def hess (t, x):  return np.sum(2*((x-t)**2 - 1)/(1+(x-t)**2)**2)

for n in (10, 15, 20):                  # Homework 2 Q4
    x = rng.standard_cauchy(n)          # theta = 0
    t = np.median(x)                    # consistent starting value
    for _ in range(100):
        step = score(t, x)/hess(t, x)
        t -= step
        if abs(step) < 1e-12: break
    print(n, np.median(x), t, x.mean())
\end{lstlisting}

\begin{lstlisting}[language=R]
set.seed(2026)
for (n in c(10,15,20)) {
  x   <- rcauchy(n)
  nll <- function(t) sum(log(1 + (x-t)^2))
  fit <- optimize(nll, interval = c(min(x), max(x)))
  cat(n, median(x), fit$minimum, mean(x), "\n")
}
\end{lstlisting}

\noindent
\textbf{Always verify against a grid or a built-in optimiser.} The likelihood can be
multimodal, and a single Newton run tells you nothing about whether you found the global
maximum.

%%%%%%%%%%%%%%%%%%%%%%%%%%%%%%%%%%%%%%%%%%%%%%%%%%%%%%%%%
\section{When Newton--Raphson fails --- and scoring does not}

This is the most instructive example in the document, and it justifies the Lecture 8
substitution far better than ``it is cheaper''.

Take four observations in two well-separated clusters:
\[
x=(-5,\ -4.5,\ 4.5,\ 5),\qquad \text{median}=0 .
\]
The log-likelihood is bimodal, with \emph{three} stationary points:
\begin{center}
\renewcommand{\arraystretch}{1.3}
\begin{tabular}{@{}cccl@{}}
\toprule
$\theta$ & $\ell(\theta)$ & $\ell''(\theta)$ & type \\
\midrule
$-4.6216$ & $-9.1205$ & $<0$ & global maximum \\
$0.0000$ & $-12.6289$ & $\mathbf{+0.3125}$ & \textbf{local MINIMUM} \\
$+4.6216$ & $-9.1205$ & $<0$ & global maximum \\
\bottomrule
\end{tabular}
\end{center}

\begin{warn}{The sample median sits exactly on the local minimum}
Newton--Raphson solves $\ell'=0$. It does \textbf{not} care whether the root is a maximum or a
minimum. Started at $\theta^{(0)}=0$ we have $\ell'(0)=0$ already, so the step is zero and the
iteration \emph{never moves} --- it reports a local minimum as the answer.

Worse, this is not a knife-edge accident. Starting at $\theta^{(0)}=4.0$, well inside the
right-hand mode:
\[
4.000\ \to\ 5.482\ \to\ 3.293\ \to\ 0.330\ \to\ 0.002\ \to\ 0 .
\]
It overshoots to $5.48$, comes back to $3.29$ where $\ell''$ has flipped sign to
\emph{positive}, and from there gets dragged into the minimum. Only a start at $4.5$ ---
already very close --- converges correctly to $4.6216$.
\end{warn}

\begin{idea}{Scoring succeeds, and here is why}
The scoring update $\theta^{(m+1)}=\theta^{(m)}+\frac{2}{n}\ell'(\theta^{(m)})$ from
$\theta^{(0)}=0.1$ converges to $4.6216$: the global maximum.

The reason is structural. Scoring divides by $-nI(\theta)$, which is \textbf{always negative},
so the step is always a positive multiple of $\ell'$ --- i.e.\ always in the direction of
\emph{increasing} likelihood. It is a gradient ascent with a principled step size and
\textbf{cannot be attracted to a minimum}. Newton--Raphson divides by $\ell''$, whose sign
flips wherever the likelihood is locally convex, and there the step points the wrong way.

\smallskip
So the Lecture 8 substitution buys more than cheapness: it buys an \textbf{ascent guarantee}.
\end{idea}

\begin{recipe}{Practical rules}
\begin{enumerate}
\item Start from the sample median (Lecture 10) --- but treat it as a suggestion, not a
guarantee. Check $\ell''<0$ there.
\item Prefer \textbf{scoring} when the likelihood may be multimodal; prefer \textbf{NR} for
speed once you are already near the maximum. Many implementations do scoring first, then
switch.
\item \textbf{Always confirm the answer is a global maximum} by evaluating $\ell$ on a grid, or
by restarting from several points.
\item If NR produces a $\theta$ with $\ell''(\theta)>0$, you have found a minimum. Discard it.
\end{enumerate}
\end{recipe}

%%%%%%%%%%%%%%%%%%%%%%%%%%%%%%%%%%%%%%%%%%%%%%%%%%%%%%%%%
\section{Beyond the Cauchy (not examinable, but worth seeing)}

The Cauchy is the course's only instance, which can leave the impression that Newton--Raphson
is a Cauchy trick. It is not. The Gamma with \emph{both} parameters unknown --- whose
exponential-family form you derived in \texttt{PI\_Exponential\_Families\_Crash\_Course.pdf}
--- is another natural case.

For $X_1,\dots,X_n\iid\mathrm{Gamma}(\alpha,\beta)$ (shape, scale),
\[
\ell(\alpha,\beta)=(\alpha-1)\sum_i\log x_i-\frac{1}{\beta}\sum_i x_i-n\log\Gamma(\alpha)-n\alpha\log\beta .
\]
Setting $\partial\ell/\partial\beta=0$ gives $\hat\beta=\Xbar/\alpha$; substituting into
$\partial\ell/\partial\alpha=0$ leaves a \emph{one}-dimensional equation,
\[
\boxed{\ \log\alpha-\psi(\alpha)=s,\qquad
s:=\log\Xbar-\frac1n\sum_i\log x_i=\log\frac{\text{AM}}{\text{GM}}\ >0,\ }
\]
where $\psi$ is the digamma function. No closed form --- so, Newton--Raphson, with
\[
g(\alpha)=\log\alpha-\psi(\alpha)-s,\qquad g'(\alpha)=\frac1\alpha-\psi'(\alpha).
\]

\textbf{Worked instance.} $x=(1.2,\,3.4,\,0.8,\,5.1,\,2.6,\,1.9,\,4.3,\,2.2)$, $n=8$:
$\Xbar=2.6875$, $\frac1n\sum\log x_i=0.8321$, so $s=0.15653$. The standard starting
approximation $\alpha_0=\dfrac{3-s+\sqrt{(s-3)^2+24s}}{12s}$ gives $3.34579$, and then

\begin{center}
\renewcommand{\arraystretch}{1.3}
\begin{tabular}{@{}ccc@{}}
\toprule
$m$ & $\alpha^{(m)}$ & $g(\alpha^{(m)})$ \\
\midrule
0 & $3.345793$ & $+2.89\times10^{-4}$ \\
1 & $3.351692$ & $+5.31\times10^{-7}$ \\
2 & $3.351703$ & $+1.80\times10^{-12}$ \\
\bottomrule
\end{tabular}
\end{center}
\[
\hat\alpha=3.3517,\qquad \hat\beta=\Xbar/\hat\alpha=0.8018 .
\]
Cross-checked against a bracketing root-finder on $[0.1,50]$: identical to six decimals.

\begin{idea}{The moral}
The MLE has no closed form here for the same underlying reason as in the Cauchy: the
likelihood equation involves a transcendental function of the parameter. Note the connection
to the exponential-family derivation --- since $T_1(x)=\log x$ for the Gamma, the natural
estimating equation is about $\E[\log X]=\psi(\alpha)+\log\beta$, and inverting $\psi$ is
exactly what has no closed form.
\end{idea}

%%%%%%%%%%%%%%%%%%%%%%%%%%%%%%%%%%%%%%%%%%%%%%%%%%%%%%%%%
\section{Related questions, with sources}

\begin{enumerate}[label=\textbf{Q\arabic*.}]
\item \textbf{[Homework 2 Q4 --- solved in \texttt{PI\_Assignment\_Solutions.pdf} \S2]}
Generate $X_1,\dots,X_n\iid\mathrm{Cauchy}(0)$ for $n=10,15,20$, compute the MLE, and compare
with $\Xbar$.

\item \textbf{[Lecture 8 --- worked in class]} Write down the method-of-scoring update for the
Cauchy location parameter and explain why $I(\theta)$ does not involve $\theta$.

\item \textbf{[Lecture 7 --- stated in class, not set]} Show the Cauchy likelihood equation is
a polynomial of degree $2n-1$, and deduce that the MLE need not be unique.

\item \textbf{[Lecture 10 --- stated in class]} Why start the iteration from the sample median
rather than the sample mean? \emph{(Because $\Xbar\sim\mathrm{Cauchy}(\theta,1)$ for every $n$,
so it is not even consistent, while the median is.)}

\item \textbf{[Class Exercise 10.2 --- solved in
\texttt{PI\_Class\_Exercises\_Solved.pdf} \S10]} Prove consistency of the MLE in general, and
explain why $U[0,\theta]$ is consistent despite failing the regularity conditions.

\item \textbf{[New --- \S6 above]} For $x=(-5,-4.5,4.5,5)$, find all stationary points of the
Cauchy log-likelihood, classify each, and show that Newton--Raphson started at the sample
median fails while the method of scoring succeeds.
\end{enumerate}

\begin{idea}{What a written exam could realistically ask}
Not ``perform five iterations by hand'' --- that is a lab task. Far more likely, and worth
roughly 8--10 marks:
\begin{quote}
\emph{Derive the likelihood equation for the Cauchy location parameter; show it has no
closed-form solution; describe the Newton--Raphson and method-of-scoring iterations; compute
$I(\theta)$ and hence give the scoring update explicitly; and state, with reason, a suitable
starting value.}
\end{quote}
Every ingredient of that answer is in \S2--\S3 of this document. That is also why
Newton--Raphson does not appear in \texttt{PI\_Mock\_Exam.pdf}.
\end{idea}

\vfill
\begin{center}
\small\itshape
Companions: \texttt{PI\_Assignment\_Solutions.pdf} (HW2 Q4),
\texttt{PI\_Class\_Exercises\_Solved.pdf} \S10,
\texttt{PI\_Exponential\_Families\_Crash\_Course.pdf},
\texttt{PI\_Exam\_Revision.pdf} \S12--\S13.
\end{center}

\end{document}
